% Corrected French for printed p.203.
% Three print-level repairs are recorded append-only.
\clearpage
\SourcePage{203}
\begingroup
\setlength{\abovedisplayskip}{6pt plus 1pt minus 1pt}
\setlength{\belowdisplayskip}{6pt plus 1pt minus 1pt}
\enlargethispage{\baselineskip}

\LayerTarget{tohoku:remark:5.2.4:2}
\noindent et celle de $G$. De façon générale, supposons données deux suites
exactes à 5 termes comme dans le cas actuel (cf. diagramme ci-dessous).

\LayerTarget{tohoku:diagram:five-term-comparison-5.2}
\begin{center}
\setlength{\unitlength}{0.62mm}
\begin{picture}(140,78)
\put(58,75){\makebox(0,0){$0$}}
\put(58,71){\vector(0,-1){6}}
\put(58,61){\makebox(0,0){$A'$}}
\put(58,57){\vector(0,-1){9}}
\put(54,52){\makebox(0,0)[r]{$\alpha$}}
\put(2,44){\makebox(0,0){$0$}}
\put(5,44){\vector(1,0){12}}
\put(22,44){\makebox(0,0){$A$}}
\put(27,44){\vector(1,0){25}}
\put(40,48){\makebox(0,0){$\alpha'$}}
\put(58,44){\makebox(0,0){$H^1$}}
\put(64,44){\vector(1,0){21}}
\put(75,39){\makebox(0,0){$\beta'$}}
\put(91,44){\makebox(0,0){$B$}}
\put(96,44){\vector(1,0){12}}
\put(102,39){\makebox(0,0){$\gamma'$}}
\put(113,44){\makebox(0,0){$C$}}
\put(118,44){\vector(1,0){11}}
\put(123,39){\makebox(0,0){$\delta'$}}
\put(135,44){\makebox(0,0){$D$}}
\put(58,40){\vector(0,-1){7}}
\put(62,36){\makebox(0,0)[l]{$\beta$}}
\put(58,29){\makebox(0,0){$B'$}}
\put(58,25){\vector(0,-1){5}}
\put(62,22){\makebox(0,0)[l]{$\gamma$}}
\put(58,17){\makebox(0,0){$C'$}}
\put(58,13){\vector(0,-1){5}}
\put(62,10){\makebox(0,0)[l]{$\delta$}}
\put(58,5){\makebox(0,0){$D'$}}
\put(62,59){\vector(2,-1){25}}
\put(76,58){\makebox(0,0){$u$}}
\put(26,42){\vector(2,-1){28}}
\put(39,31){\makebox(0,0){$u'$}}
\end{picture}
\end{center}

\LayerTarget{tohoku:construction:five-term-comparison-maps}
On a posé $u=\beta'\alpha$, $u'=\beta\alpha'$. On a
$A\cap A'=\operatorname{Ker}u=\operatorname{Ker}u'$. On a des
homomorphismes naturels
\[
\begin{aligned}
\beta\beta'^{-1}&:\operatorname{Ker}\gamma'\longrightarrow
  \operatorname{Coker}u',\\
\beta'\beta^{-1}&:\operatorname{Ker}\gamma\longrightarrow
  \operatorname{Coker}u.
\end{aligned}
\]
Ces homomorphismes donnent lieu aux deux suites exactes suivantes (où $H^1$ ne
figure plus) :

\LayerTarget{tohoku:equation:5.2.12}
\[
\tag{5.2.12}
\left\{
\begin{aligned}
0\longrightarrow{}&\operatorname{Ker}u\longrightarrow A'
\longrightarrow\operatorname{Ker}\gamma'
\longrightarrow\operatorname{Coker}u'\longrightarrow C'\longrightarrow D',\\[-0.2em]
&\mathrel{\Vert}\\[-0.45em]
0\longrightarrow{}&\operatorname{Ker}u'\longrightarrow A
\longrightarrow\operatorname{Ker}\gamma
\longrightarrow\operatorname{Coker}u\longrightarrow C\longrightarrow D.
\end{aligned}
\right.
\]

\LayerTarget{tohoku:example:five-term-comparison}
On trouve par exemple : Soit $u$ l'homomorphisme naturel de $H^1(Y,A^G)$ dans
$H^1(X,A)^G$, $u'$ l'homomorphisme naturel de $H^1(G,\Gamma(A))$ dans
$H^0(Y,\mathbf H^1(G,A))$ ; \textit{alors le noyau de $u$ est isomorphe à
celui de $u'$, tandis que l'image de $u$ est isomorphe au noyau de
l'homomorphisme naturel de
$\operatorname{Ker}(H^1(X,A)^G\longrightarrow H^2(G,\Gamma(A)))$ dans
$\operatorname{Coker}u'$.}

\LayerTarget{tohoku:remark:5.2.4:3}
\noindent 3. Les conditions de la proposition 5.2.2 sont vérifiées dans la
plupart des cas qu'on rencontre dans la nature. Son intérêt réside surtout
dans la formule (5.2.11), qui, lorsque $G$ est fini, peut aussi s'écrire

\LayerTarget{tohoku:equation:5.2.11-bis}
\[
\mathbf H^q(G,A)(y)=H^q(G,f_*(A)(y)),\qquad y\in Y.
\tag{5.2.11 bis}
\]

\LayerTarget{tohoku:application:finite-group-algebraic-variety}
et donne un calcul explicite des faisceaux $\mathbf H^q(G,A)$. La proposition
5.2.2 est applicable chaque fois que $G$ est fini et $X$ séparé. Un autre cas
important est celui où $G$ est un groupe fini d'automorphismes d'une variété
algébrique abstraite $X$, tel que $Y=X/G$ soit une variété algébrique (par
exemple $G$ est le groupe de Galois d'un revêtement $X$, ramifié ou non, d'une
variété normale $Y$) : les conditions de la proposition 5.2.2 sont alors
vérifiées si $A$ est un \textit{$G$-faisceau algébrique cohérent}, ou si $A$
est le faisceau multiplicatif $\mathbf O_X^*$ des germes de fonctions
régulières inversibles sur $X$. De plus, si $X$ est non ramifiée sur $Y$,
alors on montre que $\mathbf H^n(G,A)=0$ pour $n>0$, quand $A$ est un
$G$-faisceau algébrique cohérent, ce qui permet d'appliquer prop. 5.2.3. Dans
le même cas, on peut aussi montrer que $\mathbf H^1(G,\mathbf O_X^*)=0$. On
pourrait aussi considérer des variétés arithmétiques, les mêmes résultats sont
valables.

\LayerTarget{tohoku:section:5.3}
\bigskip
\noindent\textbf{5.3. Cas d'un groupe discontinu d'homéomorphismes.} Pour
abréger et par abus de langage, nous dirons que $G$ est un \textit{groupe
discontinu d'homéomorphismes de $X$} s'il satisfait à la condition suivante :

\LayerTarget{tohoku:condition:5.3:D}
\noindent\textit{(D) Pour tout $x\in X$, le stabilisateur $G_x$ de $x$ est
fini, et il existe un voisinage $V_x$ de $x$ tel que pour tout $g\in G$ non
dans $G_x$, on ait $g\cdot V_x\cap V_x=\phi$.}

\LayerTarget{tohoku:remark:condition-D-basic-properties:p1}
On peut alors supposer évidemment $V_x$ ouvert, et de plus
$g\cdot V_x=V_x$ pour $g\in G_x$ (en remplaçant au besoin $V_x$ par
l'intersection des $g\cdot V_x$, $g\in G_x$). La condition (D) est vérifiée
par exemple si $G$ est fini et de plus $X$ séparé, comme on voit aussitôt. Si
$X$ est une variété algébrique irréductible munie de sa

\endgroup
\TohokuContinuation{204}
